\section{Moral outrage: this should not happen in America}
\label{sec:outrage}

Many health laws gain momentum after a perceived injustice: an insurance denial,
an unaffordable treatment, a discriminatory algorithm, or unequal access to care.
The framing becomes a sentence a committee can repeat on the floor: this is
unfair, and it should not happen here. Outrage is effective because it draws both
public attention and press coverage to a specific, fixable wrong.

[DRAFTING INSTRUCTIONS: in the full-narrative, anchor the outrage appeal in
\texttt{review/references/narrative\_refs.bib} entry
\texttt{obermeyer2019dissecting} (racial bias in a widely used clinical
algorithm) so the injustice is documented, not asserted. Connect it to the
bill's transparency and bias-surveillance provisions and to
\texttt{review/references/author\_works.bib} entry
\texttt{Kawchak2026PatientPriorityProposedU} (patient priority of proposed
bills). Reproduce Mermaid figure
\texttt{review/mermaid/diagrams/15-vulnerable-population-safeguards.md} as colored
TikZ. Add a body-width table contrasting the unverified status quo with the
remedy each H. R. 9510 section provides.]

\begin{narrativefig}
\node[nfrisk] (a) at (0,0) {Unequal access};
\node[nfdec] (b) at (3.8,0) {Bias\\surveillance};
\node[nfact] (c) at (7.6,0) {Equal protection};
\draw[nfedge] (a)--(b); \draw[nfedge] (b)--(c);
\end{narrativefig}
\vspace{-0.2cm}
\figcaption{Figure 4. From an unfair gap to an audited safeguard (placeholder;
expanded from Mermaid 15 in the full-narrative).}

The remedy is not outrage for its own sake. It is a precise statutory check that
turns a documented disparity into a measured, reported, and correctable number.
